\documentclass[11pt,a4paper]{appletmanual}
\usepackage[spanish,es-noquoting,es-nodecimaldot]{babel}
\manuallang{es}

\appletname{Óptimo del Consumidor}
\appletsub{Dos bienes, preferencias monótonas y una restricción presupuestaria
lineal. El diagrama 2D y la superficie 3D son lo mismo visto de dos maneras:
mueve algo en cualquiera de los dos y el otro responde.}
\appletdir{consumer-optimum}

\begin{document}
\manualtitlepage{Manual del applet · Microeconomía I · Universidad de los Andes\\
Traducido del applet de GeoGebra «Maximización de Utilidad, Curvas de Nivel y…».\\
Código MIT. Se abre en el navegador; no hay nada que instalar.}

\manualtoc
\thispagestyle{empty}
\clearpage
\setcounter{page}{1}

% ==========================================================================
\section{Qué es esto}

El problema del consumidor se enseña en dos dimensiones: curvas de indiferencia
y una recta de presupuesto. Pero la utilidad es una función de dos variables, y
las curvas de indiferencia son sus curvas de nivel. El diagrama de siempre es
la vista desde arriba de una superficie.

Este applet muestra las dos vistas a la vez y las mantiene enganchadas. Arrastre
la canasta en el plano y se mueve en la superficie; arrástrela en la superficie
y se mueve en el plano. La proyección ortogonal X--Z convierte «encontrar la
mejor canasta asequible» en «encontrar la cima de esta loma».

\figher[1.0]{panels}{Los dos paneles. Izquierda: el diagrama de manual, con el
mapa de utilidad debajo. Derecha: la misma función como superficie
$z = u(x, y)$, con la recta presupuestaria levantada sobre ella.}

\begin{amkey}
La curva presupuestaria levantada sobre la superficie es el objeto que une las
dos vistas. Su punto más alto \emph{es} el óptimo. Verlo como una loma con una
cima quita el misterio a la condición de tangencia.
\end{amkey}

\subsection{Para quién}

Para el tema de elección del consumidor, y para cualquier curso donde haga
falta que las curvas de nivel de una función de dos variables dejen de ser un
dibujo abstracto. Trae además un modo de retos con seis familias de
preferencias, puntuación y progreso guardado, que funciona bien como tarea.

%% Secciones comunes a todos los manuales, edición española.
%% Se incluyen con \input{common-es}. Lo único que cambia entre applets es
%% \appletfolder, que es también el nombre de la carpeta y de la ruta en el sitio.

\section{Cómo abrirlo}

No hay que instalar nada. La aplicación entera es un solo archivo
\code{index.html}: no descarga librerías, no llama a ningún servidor y no
guarda nada fuera del navegador.

\subsection{Las tres maneras}

\begin{description}
\item[Doble clic sobre el archivo.] Es lo más rápido. Busque
  \code{\appletfolder/index.html} y ábralo. El navegador lo muestra desde el disco,
  con la barra de direcciones empezando por \code{file://}. Funciona sin
  conexión desde el primer momento.
\item[Desde la web.] Si alguien ya lo publicó, la dirección termina en
  \code{/\appletfolder/}. Una vez cargada la página, puede desconectarse: no
  necesita la red para nada más.
\item[Servido desde una carpeta.] Para una clase, poner la carpeta en
  cualquier servidor estático basta. Con Python instalado:
  \begin{quote}\code{python3 -m http.server 8000}\end{quote}
  y luego \code{http://localhost:8000/\appletfolder/} en el navegador.
\end{description}

\begin{amnote}
Guardar la página con \key{Ctrl}\,+\,\key{S} (o \key{Cmd}\,+\,\key{S} en
Mac) deja una copia propia que sigue funcionando sin conexión y que se puede
llevar en una memoria USB o repartir por correo.
\end{amnote}

\subsection{Idioma y tema}

Arriba a la derecha hay dos pares de botones.

\begin{ctrltable}{Control & Qué hace}
ES / EN & Cambia el idioma de toda la interfaz al vuelo. Nada se recalcula:
          puede cambiarlo en mitad de una explicación sin perder el estado. \\
Oscuro / Claro & Cambia el tema. El oscuro es el diseño original; el claro
          existe para proyectores y para imprimir. \\
\end{ctrltable}

Las dos elecciones se recuerdan en el navegador, así que la próxima vez se abre
como la dejó.

\subsection{Qué navegador}

Cualquiera que sea razonablemente reciente: Chrome, Edge, Firefox o Safari, en
ordenador, tableta o teléfono. En pantallas estrechas los paneles se apilan uno
sobre otro en lugar de ponerse al lado. No hace falta cuenta, ni permisos, ni
extensiones.


% ==========================================================================
\section{La pantalla}

\fig[1.0]{interface}{La pantalla completa.}

\begin{description}
\item[Modo.] \ui{Explorar} o \ui{Retos}, arriba a la derecha.
\item[Preferencias.] La casilla de $u(x, y)$ y seis atajos.
\item[Restricción presupuestaria.] $p_x$, $p_y$, $m$.
\item[Método gráfico.] Cuál de las tres maneras de buscar el óptimo está
  activa.
\item[Capas.] Catorce, que se describen más abajo.
\item[Dominio.] Hasta dónde llegan los ejes, con un botón para ajustarlos al
  presupuesto.
\end{description}

Debajo de los dos paneles hay una barra deslizante cuyo significado cambia con
el método elegido, y a los pies de cada panel una fila de lecturas.

\subsection{Las lecturas}

\begin{ctrltable}{Lectura & Qué es}
$x$, $y$ & Las coordenadas de la canasta $P$. \\
$u(P)$ & La utilidad en $P$. \\
Gasto & Lo gastado sobre lo disponible, $p_x x + p_y y$ frente a $m$. \\
RMS & La relación marginal de sustitución en $P$, $u_x/u_y$, en positivo. \\
$p_x/p_y$ & El precio relativo. En el óptimo interior los dos coinciden. \\
Brecha & Cuánto le falta a $P$ para ser el óptimo, y una etiqueta:
  \emph{tangencia}, \emph{casi} o \emph{lejos}. \\
u máx. & La utilidad del verdadero óptimo. Es la que hay que alcanzar. \\
\end{ctrltable}

\begin{amnote}
La RMS se presenta en positivo, como $u_x/u_y$, y se compara directamente con
$p_x/p_y$. El original de GeoGebra mostraba $-u_x/u_y$ bajo una etiqueta que
decía $u_x/u_y$, lo que hacía imposible comparar los dos números a ojo.
\end{amnote}

% ==========================================================================
\section{Los tres métodos gráficos}

\begin{description}
\item[Elegir la canasta.] $P$ se mueve por la recta presupuestaria, y el
  deslizador de abajo la recorre de $x = 0$ a $x = m/p_x$. La tarea es subir
  mientras el color bajo $P$ siga aclarándose.
\item[Subir la curva de nivel.] $P$ se queda quieta y el deslizador sube el
  nivel $u$. La tarea es encontrar el nivel en el que la curva pasa de cortar
  la recta dos veces a no cortarla: ahí la roza en un solo punto.
\item[Punto libre.] $P$ se mueve por todo el cuadrante, dentro o fuera del
  conjunto factible. Sirve para hablar de canastas inasequibles y de las que
  dejan dinero sin gastar.
\end{description}

Los dos paneles son arrastrables en los tres métodos. En 3D el puntero se
lanza como un rayo contra el campo de alturas, así que la canasta cae donde el
cursor toca visualmente el terreno.

\fig[1.0]{optimum}{Con la capa \ui{Mostrar óptimo} y la tangente encendidas.}

% ==========================================================================
\section{La cámara}

Cuatro botones sobre el panel 3D. Tres son proyecciones \textbf{ortográficas}
de verdad —sin división por perspectiva—, así que las aristas paralelas siguen
paralelas y pasos iguales en el mundo son pasos iguales en la pantalla.

\begin{ctrltable}{Botón & Proyecta sobre, y qué ejes quedan en pantalla}
X--Z & El plano $y = 0$: $x$ a la derecha, $u$ arriba. \\
X--Y & El plano $z = 0$: $x$ a la derecha, $y$ arriba. \\
Y--Z & El plano $x = 0$: $y$ a la derecha, $u$ arriba. \\
3D & Ninguno: perspectiva libre, con órbita y zoom. \\
\end{ctrltable}

\fig[0.62]{xz}{La proyección X--Z. La curva presupuestaria levantada se
proyecta en un solo arco cuya cima es el óptimo.}

\ui{X--Z} es la que hay que buscar: «encontrar la mejor canasta asequible» se
convierte en «encontrar lo alto de esta loma». \ui{X--Y} reproduce el diagrama
2D exactamente, lo que hace concreta la correspondencia entre los dos paneles.
En las dos vistas laterales una curva de indiferencia, por ser un conjunto de
nivel, se proyecta en una recta horizontal.

\fig[0.62]{xy}{La proyección X--Y. Es el diagrama 2D, dibujado por el motor 3D.}

Arrastrar o mover con las flechas fuera de una vista canónica suelta el botón
pero conserva el tipo de proyección. Arrastrar la canasta funciona en todas las
vistas: la aplicación usa solo las coordenadas que la proyección activa puede
expresar, de modo que un arrastre horizontal en \ui{Y--Z} mueve a lo largo de
$y$ y no del eje $x$ colapsado.

% ==========================================================================
\section{Las capas}

\begin{ctrltable}{Capa & Qué dibuja}
Mapa de utilidad & El campo $u$ bajo una rampa de color. \\
Curvas de fondo & Curvas de nivel a intervalos regulares. \\
Curva de nivel & La curva del nivel que marca el deslizador. \\
Curva por P & La curva de indiferencia que pasa por la canasta. \\
Conjunto factible & Vela lo que no se puede pagar. \\
Corte 3D & La recta presupuestaria levantada sobre la superficie. \\
Tangente en P & La tangente a la curva de indiferencia en $P$. \\
Gradiente & $\nabla u$ en $P$. \\
Parciales & $u_x$ y $u_y$ por separado. \\
Plano de corte & El plano vertical sobre la recta presupuestaria. \\
Retícula base & La cuadrícula del suelo en 3D. \\
Sólo lo asequible & Recorta la superficie al conjunto factible. \\
Mostrar óptimo & Revela el verdadero óptimo. \\
Dejar rastro & Guarda las curvas de nivel que se van barriendo. \\
\end{ctrltable}

% ==========================================================================
\section{Casos particulares}

Los seis atajos de preferencias cubren los casos que un curso necesita. Dos
merecen mirarse despacio porque rompen la condición de tangencia.

\fig[1.0]{complements}{Complementarios perfectos, \code{min(2x,y)}. La curva de
indiferencia tiene un codo y no hay tangente que igualar.}

\fig[1.0]{substitutes}{Sustitutos perfectos, \code{2x+3y}. Con
$p_x/p_y = 2/3$ la recta presupuestaria y las curvas de indiferencia son
paralelas: toda la recta es óptima.}

\begin{amnote}
Con los precios de partida ($p_x = 2$, $p_y = 3$) el atajo \code{2x+3y} da
exactamente $\text{RMS} = p_x/p_y$. Es el caso degenerado en que el consumidor
es indiferente entre todas las canastas asequibles, y merece la pena enseñarlo
antes de que aparezca por sorpresa en un examen.
\end{amnote}

% ==========================================================================
\section{El modo de retos}

\fig[1.0]{challenge}{El modo de retos.}

Seis niveles, cada uno con su familia de preferencias y con precios y renta
generados de nuevo en cada intento:

\begin{ctrltable}{Nivel & Familia y tarea}
1 & Cobb--Douglas. Colocar $P$ en la canasta óptima. \\
2 & Cobb--Douglas asimétrica. Subir la curva de nivel hasta que roce. \\
3 & Sustitutos perfectos. Colocar $P$. \\
4 & Complementarios perfectos. Colocar $P$. \\
5 & Cuasilineal. Subir la curva de nivel. \\
6 & CES. Colocar $P$. \\
\end{ctrltable}

Mientras un reto está vivo, los precios, la renta y la función de utilidad
quedan bloqueados. Comprobar la respuesta revela el óptimo, puntúa el intento y
explica \emph{cuál} condición de optimalidad se aplicaba: tangencia, vértice o
esquina. El progreso se guarda en el navegador.

\begin{amnote}
Que la tarea alterne entre «colocar la canasta» y «subir la curva» es
deliberado: son los dos métodos gráficos del temario, y un estudiante que solo
practique uno no reconoce el otro en un examen.
\end{amnote}

% ==========================================================================
\section{Actividades para clase}

\begin{activity}{Las dos vistas son una}
\amset Valores de partida ($u = x^{0,5}y^{0,5}$, $p_x = 2$, $p_y = 3$,
$m = 24$). Método \ui{Elegir la canasta}.
\amexp Arrastre $P$ en el panel 2D y mírela moverse en el 3D. Después
arrástrela en el 3D y mírela moverse en el 2D.
\par\smallskip
Pase la cámara a \ui{X--Y}: el panel derecho se convierte en el izquierdo. Es la
demostración más corta de que una curva de indiferencia es una curva de nivel.
\end{activity}

\begin{activity}{Encontrar la cima}
\amset Cámara en \ui{X--Z}, método \ui{Elegir la canasta}.
\amexp La curva presupuestaria levantada se proyecta en un arco. Mueva el
deslizador y vea $u(P)$ subir hasta un máximo y volver a bajar.
\par\smallskip
La lectura \ui{u máx.} dice $4{,}899$, que es $\sqrt{24}$. El óptimo está en
$x^\ast = 6$, $y^\ast = 4$. Compruébelo con la fórmula de Cobb--Douglas,
$x^\ast = \alpha m/p_x$, y vea que el número coincide con la cima del arco.
\end{activity}

\begin{activity}{La tangencia, encontrada de dos maneras}
\amset Método \ui{Elegir la canasta}; después \ui{Subir la curva de nivel}.
\amexp Con el primer método, suba por la recta mientras el color bajo $P$ se
siga aclarando: la lectura \ui{Brecha} baja hasta decir \emph{tangencia}.
\par\smallskip
Con el segundo, suba $u$ hasta que la curva pase de cortar la recta dos veces a
no cortarla. El nivel en el que la roza es el mismo $4{,}899$. Los dos métodos
del temario dan el mismo número, y conviene que la clase lo vea en el mismo
minuto.
\end{activity}

\begin{activity}{RMS contra precio relativo}
\amset Valores de partida, capa \ui{Tangente en P} encendida.
\amexp Con $P$ lejos del óptimo, \ui{RMS} $= 1{,}33$ y $p_x/p_y = 0{,}67$: la
tangente es más inclinada que la recta presupuestaria. Mueva $P$ hasta que los
dos números coincidan.
\par\smallskip
Ahí está la condición de primer orden, leída como dos números que se
persiguen. Encienda \ui{Gradiente} y \ui{Parciales} para ver de dónde sale la
RMS.
\end{activity}

\begin{activity}{Un óptimo sin tangencia}
\amset Atajo \ui{Complementarios / Complements}, \code{min(2x,y)}.
\amexp La lectura \ui{u máx.} da $6{,}000$, en $x^\ast = 3$, $y^\ast = 6$. Pero
la RMS en el codo no está definida: la lectura da $0{,}00$ y ninguna tangente
iguala nada.
\par\smallskip
Pregunte cómo se encuentra entonces el óptimo. La respuesta es que se resuelve
el sistema $2x = y$ junto con la recta presupuestaria, y que la condición de
tangencia no interviene. Es el caso que rompe la receta.
\end{activity}

\begin{activity}{Todas las canastas empatan}
\amset Atajo \ui{Sustitutos / Substitutes}, \code{2x+3y}, precios de partida.
\amexp $\text{RMS} = 2/3$ y $p_x/p_y = 2/3$: iguales. La recta presupuestaria y
las curvas de indiferencia son paralelas, así que \emph{cualquier} canasta de
la recta es óptima.
\par\smallskip
Mueva $p_x$ un poco arriba o abajo y el óptimo salta de golpe a un eje. Es la
demanda de esquina, y este es el único ajuste en que se la puede ver justo en
el punto de bascular.
\end{activity}

\begin{activity}{Lo que no se puede pagar}
\amset Método \ui{Punto libre}. Capa \ui{Conjunto factible} encendida.
\amexp Arrastre $P$ fuera del velo. La lectura \ui{Gasto} pasa de $24{,}0$ y
el applet no lo impide: la canasta existe, simplemente no se puede pagar.
\par\smallskip
Arrástrela después hacia adentro, dejando dinero sin gastar. Con preferencias
monótonas eso nunca es óptimo, y se ve porque siempre hay una canasta al
noreste con más color.
\end{activity}

\begin{activity}{Retos como tarea}
\amset Modo \ui{Retos}.
\amexp Los seis niveles generan precios y renta nuevos en cada intento, así que
no se puede memorizar la respuesta. Al comprobar, el applet explica qué
condición se aplicaba en ese caso.
\par\smallskip
Es una tarea de una sesión con corrección automática. El progreso se guarda en
el navegador del estudiante, así que conviene pedir una captura de pantalla de
la puntuación.
\end{activity}

% ==========================================================================
\section{Sintaxis de expresiones}

\begin{ctrltable}{Elemento & Qué se admite}
Operadores & \code{+ - * / \^{} ( )} \\
Funciones & \code{sqrt ln log log10 exp abs min max sin cos tan} \\
Constantes & \code{e}, \code{pi} \\
Multiplicación implícita & \code{2xy} es \code{2*x*y}; \code{x\^{}(1/3)y\^{}(2/3)}
  se lee como se espera \\
Asociatividad & \code{\^{}} liga por la derecha y más fuerte que el menos
  unario: \code{-x\^{}2} es $-(x^2)$ \\
\end{ctrltable}

% ==========================================================================
\section{Cómo está hecho}

Todo está escrito a mano para que la página siga siendo autónoma.

\begin{ctrltable}{Pieza & Enfoque}
Lenguaje de expresiones & Tokenizador $\rightarrow$ analizador de Pratt
  $\rightarrow$ árbol de cierres. Sin \code{eval} ni \code{new Function}, así
  que una política de seguridad estricta no lo rompe. \\
Derivadas parciales & Derivación simbólica sobre el árbol, con plegado de
  constantes. Pasa a diferencias centradas cuando aparecen \code{min},
  \code{max} o \code{abs} —justo los casos con quiebre. \\
Curvas de nivel & Marching squares sobre un campo muestreado de
  $141\times141$, con desambiguación de sillas y salto de celdas indefinidas. \\
Óptimo & Barrido denso y refinamiento por sección áurea sobre la recta
  presupuestaria, y después clasificación en interior, vértice o esquina.
  Robusto donde una condición de primer orden no encuentra nada. \\
Superficie 3D & Algoritmo del pintor sobre canvas 2D. Un campo de alturas se
  ordena exactamente por profundidad de centroide, así que no hace falta
  búfer de profundidad ni hay modo de fallo por pérdida de contexto. \\
Selección en 3D & Se lanza el rayo del puntero contra el campo de alturas y se
  biseca. Los cruces cuentan en las dos direcciones, porque en una proyección
  lateral el rayo va a ras del terreno y puede emerger desde debajo. \\
\end{ctrltable}

% ==========================================================================
\section{Diferencias con el original de GeoGebra}

\begin{description}
\item[Rampa de color.] El original teñía las curvas de nivel con tres funciones
  lineales a trozos (\code{redfun}, \code{greenfun}, \code{bluefun}). El
  reemplazo va de azul y verde abajo a amarillo y naranja en medio, y luego a
  rojo vivo y marrón oscuro arriba. Como el extremo alto se oscurece, la altura
  no se recupera solo del brillo: eso lo llevan la barra de color y las
  lecturas, y toda curva dibujada sobre el mapa lleva un contorno blanco para
  seguir siendo legible sobre los marrones.
\item[Marcas.] La canasta, la recta, la tangente, el gradiente, las parciales y
  el óptimo revelado son colores fijos y vivos con contorno oscuro, ninguno de
  ellos repitiendo un tono por el que pase la rampa.
\item[Rango de color.] \code{uMIN}/\code{uMAX} se leían en dos esquinas del
  dominio. Para $\ln(x) + y$ esa esquina inferior es $-\infty$ y la rampa se
  colapsaba. La reescritura toma el percentil 2 del campo muestreado.
\item[Dominio.] Los deslizadores $x_C$, $y_C$ llegaban a $-10$, lo que deja
  \code{sqrt(x)} y \code{ln(x)} fuera de su dominio. El cuadrante está ahora
  recortado a $x, y \ge 0$.
\item[Parametrización de la recta.] $x_s$ recorría $[x_C, x_G]$ en vez de
  $[0, m/p_x]$, así que la «canasta» podía quedar en $x$ negativo con $y$ por
  encima de la recta. Ahora se parametriza por $t \in [0, m/p_x]$.
\item[Precios.] $p_X$ y $p_Y$ empezaban en 0, con lo que
  $y_s = (m - p_X x)/p_Y$ y el dominio $m/p_X$ dividían por cero. Los dos están
  acotados por debajo en $0{,}2$.
\item[Etiquetado de la RMS.] El original mostraba $-f_x/f_y$ bajo una etiqueta
  que anunciaba $u_x/u_y$: el signo no cuadraba con el nombre. Ahora se informa
  en positivo, como $f_x/f_y$, y se compara directamente con $p_x/p_y$, con la
  brecha de tangencia como lectura propia.
\item[Objetos muertos.] \code{afin}, \code{text1} (que se refería a
  $a_0, a_1, a_2$, inexistentes en el archivo), \code{BooleFuncion},
  \code{BooleTransparente}, \code{BooleMostrarVar},
  \code{BooleIntersecciones}, \code{BooleTangent}, \code{BooleTanPlane},
  \code{s}, \code{Impl}, \code{Boole} y \code{DX} eran restos del applet del que
  se copió. No se trasladaron.
\end{description}

% ==========================================================================
\section{Accesibilidad}

Los dos lienzos se manejan con teclado: las flechas mueven la canasta o el
nivel, y \key{+} y \key{-} acercan y alejan la vista 3D. Los temas claro y
oscuro están diseñados por separado, y los dos responden al ajuste del sistema.
Los números son de ancho tabular. El movimiento respeta
\code{prefers-reduced-motion}.

% ==========================================================================
\section{Problemas frecuentes}

\begin{ctrltable}{Síntoma & Qué pasa}
El mapa sale de un solo color & La función tiene un $-\infty$ en el dominio
  dibujado. Suba el mínimo del dominio o cambie la expresión. \\
No puedo arrastrar en 3D & Compruebe que no está en modo \ui{Retos} con los
  controles bloqueados. Fuera de eso, el arrastre funciona en las cuatro
  vistas. \\
La superficie sale a trozos & La función no está definida en parte del
  cuadrante. Es correcto: ahí no hay suelo. \\
El botón de cámara se apaga solo & Ha arrastrado fuera de la vista canónica. El
  tipo de proyección se conserva; vuelva a pulsarlo para recentrar. \\
Perdí la puntuación de los retos & Se guarda en el navegador. En ventana
  privada, o tras borrar los datos del sitio, empieza de cero. \\
\end{ctrltable}

%% Secciones comunes a todos los manuales, edición española.
%% Se incluyen con \input{common-es}. Lo único que cambia entre applets es
%% \appletfolder, que es también el nombre de la carpeta y de la ruta en el sitio.

\section{Publicarlo para una clase}

Sirve cualquier alojamiento estático: GitHub Pages, Netlify, un directorio web
de la universidad, o incluso una carpeta compartida servida por HTTP. Lo único
que hay que subir es el archivo \code{index.html} dentro de una carpeta con el
nombre que quiera que aparezca en la dirección.

El repositorio trae un flujo de trabajo de GitHub Actions
(\code{.github/workflows/pages.yml}) que publica todos los applets a la vez en
cada empuje a la rama principal. Para activarlo, en el repositorio:
\emph{Settings\,$\rightarrow$\,Pages\,$\rightarrow$\,Source: GitHub Actions}.
Es lo único que no puede hacer el propio flujo, porque crear un sitio de Pages
requiere permisos de administración que un \code{GITHUB\_TOKEN} nunca recibe.

\begin{amnote}
Para repartir el applet a estudiantes sin conexión fiable, mándeles el archivo
\code{index.html} directamente. Pesa menos que una fotografía y se abre con
doble clic.
\end{amnote}

\section{Licencia y procedencia}

El código de la aplicación es MIT. Puede usarlo, modificarlo y repartirlo, en
clase o fuera de ella, citando la procedencia.

Este applet es una reescritura autónoma de un applet original de GeoGebra. La
reescritura no es una traducción literal: donde el original tenía un error de
construcción, una división por cero alcanzable con los deslizadores, o un objeto
muerto heredado del applet del que se copió, esta versión lo corrige y lo
documenta. La sección \emph{Diferencias con el original de GeoGebra} enumera
esos cambios uno por uno.

Todo está escrito a mano y sin dependencias: el analizador de expresiones, la
derivación simbólica, el trazado de curvas de nivel y el dibujo. En particular
no se usa \code{eval}, de modo que lo que escribe un estudiante en una casilla
de texto no puede convertirse en código, y una política de seguridad de
contenidos estricta no rompe la página.


\end{document}
